\section{如何评估：pass@k、10@k 与真实排名}

\subsection{搜索覆盖与提交选择要分开看}

\videofigure{figures/pass-k-metrics.jpg}{pass@$k$ 衡量候选覆盖，10@$k$ 衡量从 $k$ 个样本中选出十个提交后的实际成功率。}{00:14:18--00:16:26}

若 $n$ 个样本中有 $c$ 个正确程序，无放回选择 $k$ 个时，常用估计为：

$$
\operatorname{pass@}k=1-\frac{\binom{n-c}{k}}{\binom{n}{k}}.
$$

\begin{itemize}
    \item $n$：总候选数；
    \item $c$：其中通过隐藏测试的候选数；
    \item $k$：允许观察或抽取的候选预算；
    \item pass@$k$：至少存在一个正确候选的概率。
\end{itemize}

10@$k$ 更接近真实竞赛：先生成 $k$ 个候选，再通过过滤与聚类选出至多十个提交。它同时考察 generator 与 selector，而不仅是“正确答案是否曾出现”。

\videofigure{figures/codecontests-results.jpg}{CodeContests 结果显示模型规模和聚类都能提高有限提交预算下的 solve rate。}{00:16:26--00:18:08}

\begin{warningbox}{pass@$k$ 高不代表系统可部署}
如果正确候选藏在几十万条程序中，但选择器无法识别，它对真实用户没有价值。部署指标必须显式包含提交预算、执行成本、延迟和误选率。
\end{warningbox}

\subsection{真实 Codeforces 评估}

\videofigure{figures/codeforces-evaluation.jpg}{AlphaCode 在十场大型 Codeforces 竞赛中现场生成、筛选并提交程序，平均排名约在参赛者中位附近。}{00:09:04--00:13:58}

课程强调这类评估比静态 benchmark 更可信：题目是在模型训练完成后出现的，系统按真实提交规则运行，并与数千名人类参赛者比较。AlphaCode 在当时达到约前 54\% 的排名，证明完整端到端程序生成不再只是玩具任务。

\videofigure{figures/alphacode-tradeoffs.jpg}{AlphaCode 的强项是完整问题求解与可重复验证，弱点是高计算成本、隐藏测试依赖和选择瓶颈。}{00:22:42--00:24:34}

\subsection{本章小结}

评估搜索系统时必须拆开三件事：模型能否生成正确候选、选择器能否在小预算下找到它、最终程序能否在真实平台通过并具备可接受效率。

